\newpage
\phantomsection
\label{prf:rational-fraction-multiplication}

\begin{remark*}[Return]
\hyperref[thm:rational-fraction-multiplication]{Return to Theorem}
\end{remark*}

\begin{theorem*}[Fraction Multiplication Rule in $\mathbb{Q}$]
For every $a,b,c,d\in\mathbb{Q}$ with $b\neq 0$ and $d\neq 0$,
\[
\frac{a}{b}\cdot \frac{c}{d}
=
\frac{a\cdot c}{b\cdot d}.
\]
\end{theorem*}

\begin{proof}
\textbf{Professional Standard Proof.}~\\
Set
\[
x:=\frac{a}{b}\cdot\frac{c}{d}.
\]
We show that $x$ satisfies the defining equation for
\[
\frac{a\cdot c}{b\cdot d}.
\]
Since $b\neq 0$ and $d\neq 0$, the product $bd$ is nonzero in the field
$\mathbb{Q}$.  Multiply by $bd$:
\[
(bd)x=(bd)\left(\frac{a}{b}\cdot\frac{c}{d}\right).
\]
By associativity and commutativity of multiplication,
\[
(bd)\left(\frac{a}{b}\cdot\frac{c}{d}\right)
=
\left(b\frac{a}{b}\right)\left(d\frac{c}{d}\right).
\]
By the definition of division,
\[
b\frac{a}{b}=a
\qquad\text{and}\qquad
d\frac{c}{d}=c.
\]
Therefore
\[
(bd)x=a\cdot c.
\]
By the definition of division in $\mathbb{Q}$, the unique rational number
whose product with $bd$ equals $a\cdot c$ is
\[
\frac{a\cdot c}{b\cdot d}.
\]
Hence
\[
x=\frac{a\cdot c}{b\cdot d},
\]
that is,
\[
\frac{a}{b}\cdot\frac{c}{d}
=
\frac{a\cdot c}{b\cdot d}.
\]
\end{proof}

\begin{remark*}[Proof structure]
The proof follows the same pattern as the addition and subtraction rules.
Rather than manipulating fraction notation directly, it introduces the
left-hand side as a temporary variable $x$ and proves that multiplying $x$
by the common denominator $bd$ yields the numerator $ac$.  Once that
equation is established, the definition of division identifies $x$
uniquely as $(ac)/(bd)$.  The whole argument is thus grounded in the field
structure of $\mathbb{Q}$ and the defining property of quotients.
\end{remark*}

\begin{remark*}[Dependencies]~\\
\begin{itemize}
  \item \hyperref[def:rational-division]{Definition~\ref*{def:rational-division}}
  \item \hyperref[thm:q-is-field]{Theorem~\ref*{thm:q-is-field}}
\end{itemize}
\end{remark*}
